\documentclass[11pt]{amsart}

\usepackage[T1]{fontenc}
\usepackage{lmodern}
\usepackage{microtype}
\usepackage{mathtools,amssymb,amsthm}
\usepackage{array}
\usepackage[hidelinks]{hyperref}

\hypersetup{
  pdftitle={An Order-128 2-Group whose Squares Form a Non-Powerful Subgroup}
}

\newtheorem{theorem}{Theorem}
\newtheorem{lemma}[theorem]{Lemma}
\newtheorem{corollary}[theorem]{Corollary}
\theoremstyle{remark}
\newtheorem*{remark}{Remark}

\DeclareMathOperator{\Syl}{Syl}

\title[A 2-group whose squares are not powerful]
{An Order-128 2-Group whose Squares Form a Non-Powerful Subgroup}

\date{}

\subjclass[2020]{20D15, 20E22, 20F05}
\keywords{powerful $p$-group, power subgroup, set of squares, Kourovka Notebook,
Sylow subgroup of the symmetric group, counterexample}

\begin{document}

\begin{abstract}
Problem 21.137 of the Kourovka Notebook, proposed by L.~E. Wilson, asks: if the
$p$\nobreakdash-th powers in a finite $p$-group form a subgroup, must that subgroup be
powerful? Its $p=2$ form asks whether, in a $2$-group of exponent~$8$ whose squares form a
subgroup, that subgroup must be abelian. We answer both in the negative with
$G=\Syl_2(S_8)=(C_2\wr C_2)\wr C_2$ of order $2^7=128$. The \emph{set} $S=\{g^2:g\in G\}$ is a subgroup of order $16$, isomorphic to $C_2\times D_8$
and hence non-abelian; since $\exp S=4$ forces $S^4=1$ while $|S'|=2$, $S$ is not powerful
in the published convention $H'\le P_2(H)=\langle h^4\rangle$ for $p=2$. The group is
exhibited both as seven explicit permutations of $\{1,\dots,8\}$ and by a $28$-relator
presentation. The odd-prime question in the second
sentence of 21.137 is a different question and is not addressed here.
\end{abstract}

\maketitle

\section{The question}

The Kourovka Notebook \cite{Kourovka} records, as Problem~21.137 with proposer
L.~E. Wilson, the following. All three of its sentences are quoted verbatim.

\begin{quote}
\textbf{21.137.} If the $p$-th powers in a finite $p$-group form a subgroup, must that
subgroup be powerful? That is, for $p\ne2$, if the $p$-th powers in a $p$-group of exponent
$p^2$ form a subgroup, must that subgroup be abelian? For a $2$-group of exponent $8$, if
the squares form a subgroup, must that subgroup be abelian? \emph{---~L.~Wilson}
\end{quote}

Here ``the $p$-th powers form a subgroup'' means that the \emph{set}
$\{g^p:g\in G\}$ is closed under multiplication, so that it coincides with the power
subgroup $P_1(G)=\langle g^p:g\in G\rangle$ it generates. One direction of the implication
21.137 asks about is a theorem: for odd $p$, in a powerful $p$-group the set of $p^k$-th
powers \emph{is} the subgroup it generates (Lubotzky and Mann \cite{LubotzkyMann} for $k=1$,
and the proposer \cite[Thm.~3.1]{Wilson2002} for every $k$). Problem 21.137 asks for the
converse.

The one definition the problem does not restate is the one that decides it. Following
Lubotzky and Mann \cite[\S1]{LubotzkyMann}, a finite $p$-group $H$ is \emph{powerful} when
\begin{equation}\label{eq:powerful}
\begin{aligned}
  H'&\le P_1(H)=\langle h^p:h\in H\rangle && \text{for odd }p,\\
  H'&\le P_2(H)=\langle h^4:h\in H\rangle && \text{for }p=2 .
\end{aligned}
\end{equation}
The shift to \emph{fourth} powers at $p=2$ is not a convenience: with $\langle
h^2\rangle$ in place of $\langle h^4\rangle$ the condition would hold in every group
whatsoever (Lemma~\ref{lem:vacuous} below), so the $p=2$ half of the headline question
would be trivially affirmative.

Convention \eqref{eq:powerful} is what makes the problem's own bridge sentence
(``That is,~\dots'') correct at $p=2$ as well as at odd $p$. If $\exp G=8$ and
$S=\{g^2:g\in G\}$ is a subgroup, then $\exp S=4$, so $S^4:=P_2(S)=1$ and
\eqref{eq:powerful} reduces to $S'=1$. Thus for $p=2$, ``powerful'' and ``abelian'' say the
same thing about $S$, and the third sentence of 21.137 is a faithful specialisation of the
first, not a weaker or a stronger question.

\begin{theorem}\label{thm:main}
Let $G=\Syl_2(S_8)=(C_2\wr C_2)\wr C_2$, a group of order $2^7=128$, and let
$S=\{g^2:g\in G\}$. Then
\begin{enumerate}
\item[(i)] $\exp G=8$;
\item[(ii)] $S$ is a subgroup of $G$, of order $16$, and $S=P_1(G)=G'=\Phi(G)$;
\item[(iii)] $S\cong C_2\times D_8$; in particular $S$ is \textbf{not abelian}, and
      $|S'|=2$;
\item[(iv)] $\exp S=4$, so $P_2(S)=1$, and therefore $S'\not\le P_2(S)$: the subgroup $S$
      is \textbf{not powerful}.
\end{enumerate}
Consequently the answer to the headline question of 21.137 is \emph{no}, already at $p=2$,
and the answer to its third sentence is \emph{no}.
\end{theorem}

\section{The object}

Write $[x,y]=x^{-1}y^{-1}xy$, so that $xy=yx\,[x,y]$. Let $G\le S_8$ be generated by the
seven permutations
\begin{center}
\begin{tabular}{ll@{\qquad}ll}
$f_1$ & $=(1\;5)(2\;6)(3\;7)(4\;8)$ & $f_5$ & $=(1\;3)(2\;4)(5\;7)(6\;8)$\\
$f_2$ & $=(3\;4)$                   & $f_6$ & $=(5\;6)(7\;8)$\\
$f_3$ & $=(5\;7)(6\;8)$             & $f_7$ & $=(1\;2)(3\;4)(5\;6)(7\;8)$\\
$f_4$ & $=(3\;4)(7\;8)$             &       &
\end{tabular}
\end{center}
These seven satisfy, and are presented by, the $28$ relators
\begin{equation}\label{eq:pres}
\begin{gathered}
 f_1^2=f_2^2=f_3^2=f_4^2=f_5^2=f_6^2=f_7^2=1,\\
 [f_2,f_1]=f_4,\quad [f_3,f_1]=f_5,\quad [f_6,f_1]=f_7,\\
 [f_5,f_2]=f_6f_7,\quad [f_4,f_3]=f_6,\quad [f_5,f_4]=f_7,\\
 \text{and }[f_a,f_b]=1\text{ for each of the remaining }15\text{ pairs }a>b.
\end{gathered}
\end{equation}
That is seven squares together with all $\binom72=21$ commutators, six of them non-trivial.
Every right-hand side in \eqref{eq:pres} is a word in generators of index strictly larger
than both entries of its commutator, so \eqref{eq:pres} is a polycyclic presentation with
all seven relative orders equal to $2$: collecting an arbitrary word to the left terminates,
and the presented group is spanned by the $128$ normal words
\begin{equation}\label{eq:nf}
 f_1^{a_1}f_2^{a_2}\cdots f_7^{a_7},\qquad a_i\in\{0,1\}.
\end{equation}
In $G$ those $128$ words are pairwise distinct, and $|G|=128$ by
Lemma~\ref{lem:structure}(b); hence \eqref{eq:pres} presents $G$ exactly, the words
\eqref{eq:nf} \emph{are} its elements. Two further remarks fix its identity. First, $G\le S_8$ and $8!=40320=128\cdot315$ with
$315$ odd, so a subgroup of $S_8$ of order $128$ is a full Sylow $2$-subgroup: $G\in
\Syl_2(S_8)$. Second, $G$ is also the closure of the seven level generators
$(1\;2)$, $(3\;4)$, $(5\;6)$, $(7\;8)$, $(1\;3)(2\;4)$, $(5\;7)(6\;8)$,
$(1\;5)(2\;6)(3\;7)(4\;8)$ of the automorphism group of the complete binary tree with
leaves $1,\dots,8$, which is $(C_2\wr C_2)\wr C_2$.

Put
\[
 N=\langle f_4,f_5,f_6,f_7\rangle .
\]

\begin{lemma}\label{lem:structure}
\begin{enumerate}
\item[(a)] $N\cong D_8\times C_2$: it has order $16$, $N'=\langle f_7\rangle$ has order
      $2$, and $\exp N=4$.
\item[(b)] $N\trianglelefteq G$ and $G/N\cong C_2^3$; in particular $|G|=128$ and
      $g^2\in N$ for every $g\in G$.
\item[(c)] Every element of $N$ is a square in $G$. Hence $S=\{g^2:g\in G\}=N$.
\item[(d)] $\exp G=8$.
\end{enumerate}
\end{lemma}

\begin{proof}
(a) Among the six pairs drawn from $f_4,f_5,f_6,f_7$, the only non-trivial commutator in
\eqref{eq:pres} is $[f_5,f_4]=f_7$. So $f_6$ and $f_7$ are central in $N$. As all four
generators are involutions, $[f_5,f_4]=f_5f_4f_5f_4=(f_5f_4)^2$, whence $(f_5f_4)^2=f_7$
and, taking inverses, $r:=f_4f_5$ satisfies $r^2=f_7$ and $r^4=f_7^2=1$. Since $f_7\ne1$,
$r$ has order exactly~$4$; and $f_4\,r\,f_4=f_4f_4f_5f_4=f_5f_4=r^{-1}$ with $f_4\notin
\langle r\rangle=\{1,r,f_7,rf_7\}$, so $\langle f_4,f_5\rangle=\langle r,f_4\rangle\cong
D_8$ of order $8$, with derived subgroup $\langle f_7\rangle$. Finally $f_6$ is central in
$N$ and, by uniqueness of the normal form \eqref{eq:nf}, $f_6\notin\langle f_4,f_5\rangle$;
so $N=\langle f_4,f_5\rangle\times\langle f_6\rangle\cong D_8\times C_2$, of order $16$ and
exponent~$4$, with $N'=\langle f_7\rangle$ of order $2$.

(b) Each of $[f_4,f_i],[f_5,f_i],[f_6,f_i],[f_7,f_i]$ for $i=1,2,3$ lies in $N$: by
\eqref{eq:pres} the only non-trivial ones are $[f_6,f_1]=f_7$, $[f_5,f_2]=f_6f_7$ and
$[f_4,f_3]=f_6$. As $N$ is generated by $f_4,\dots,f_7$ and $G$ by $f_1,\dots,f_7$, this
gives $N\trianglelefteq G$. Modulo $N$ the images of $f_1,f_2,f_3$ are involutions which
commute, because $[f_2,f_1]=f_4\in N$, $[f_3,f_1]=f_5\in N$ and $[f_3,f_2]=1$; so $G/N$ is
generated by three commuting involutions, and it has order $8$ because the $128$ normal
words \eqref{eq:nf} are distinct in $G$ while $|N|=16$. Thus $G/N\cong C_2^3$ and
$|G|=16\cdot8=128$. Since $G/N$ has exponent $2$, $(gN)^2=N$, i.e.\ $g^2\in N$, for every
$g\in G$.

(c) Table~\ref{tab:roots} lists, for each of the sixteen elements $s$ of $N$, an element
$r\in G$ with $r^2=s$. Two entries are immediate by hand: from $[f_2,f_1]=f_4$ and
$f_1^2=f_2^2=1$ we get $(f_2f_1)^2=f_2f_1f_2f_1=[f_2,f_1]=f_4$, hence
$(f_1f_2)^2=\bigl((f_2f_1)^2\bigr)^{-1}=f_4$, and likewise $(f_1f_3)^2=f_5$. Together with
(b), $S=N$.

(d) By (a) and (c), $S=N$ has exponent $4$, so $g^8=(g^2)^4=1$ for every $g\in G$ and
$\exp G$ divides $8$. Conversely $f_4f_5\in N$ has order $4$ by (a) and is a square by (c)
--- Table~\ref{tab:roots} gives $(f_1f_2f_3f_5)^2=f_4f_5$ --- so $f_1f_2f_3f_5$ has order
$8$. Hence $\exp G=8$ exactly.
\end{proof}

\begin{table}[htbp]
\caption{The sixteen elements $s$ of $S=N$, each with a square root $r\in G$ and its order.
Words are read left to right in $f_1,\dots,f_7$.}
\label{tab:roots}
\begin{tabular}{lcl@{\qquad\qquad}lcl}
\hline
$s$ & $r$ with $r^2=s$ & $|r|$ & $s$ & $r$ with $r^2=s$ & $|r|$\\
\hline
$1$        & $1$           & $1$ & $f_5f_6$      & $f_1f_3f_4$    & $4$\\
$f_4$      & $f_1f_2$      & $4$ & $f_5f_7$      & $f_1f_3f_6$    & $4$\\
$f_5$      & $f_1f_3$      & $4$ & $f_6f_7$      & $f_2f_5$       & $4$\\
$f_6$      & $f_3f_4$      & $4$ & $f_4f_5f_6$   & $f_1f_2f_3f_6$ & $8$\\
$f_7$      & $f_1f_6$      & $4$ & $f_4f_5f_7$   & $f_1f_2f_3f_4$ & $8$\\
$f_4f_5$   & $f_1f_2f_3f_5$& $8$ & $f_4f_6f_7$   & $f_1f_2f_5$    & $4$\\
$f_4f_6$   & $f_1f_2f_4f_5$& $4$ & $f_5f_6f_7$   & $f_1f_3f_4f_5$ & $4$\\
$f_4f_7$   & $f_1f_2f_6$   & $4$ & $f_4f_5f_6f_7$& $f_1f_2f_3$    & $8$\\
\hline
\end{tabular}
\end{table}

\begin{proof}[Proof of Theorem~\ref{thm:main}]
Part (i) is Lemma~\ref{lem:structure}(d). By Lemma~\ref{lem:structure}(b,c), $S=N$ is a
subgroup of order $16$, and it equals $P_1(G)=\langle g^2:g\in G\rangle$ because the set of
squares already is a subgroup. For $S=G'$: by
Lemma~\ref{lem:structure}(b) the quotient $G/N$ is abelian, so $G'\le N$, while
\eqref{eq:pres} exhibits all four generators of $N$ as commutators of elements of $G$,
namely $f_4=[f_2,f_1]$, $f_5=[f_3,f_1]$, $f_6=[f_4,f_3]$ and $f_7=[f_6,f_1]$, so $N\le G'$;
hence $G'=N=S$. Since $\Phi(G)=G'P_1(G)$ for a $p$-group, $\Phi(G)=S$ as well, and (ii)
holds. Part
(iii) is Lemma~\ref{lem:structure}(a); explicitly, $f_4=(f_1f_2)^2$ and $f_5=(f_1f_3)^2$
are two squares with
\[
 [f_4,f_5]=f_7\ne1 ,
\]
so $S$ is not abelian, and $S'=\langle f_7\rangle$ has order $2$. For (iv), $\exp S=4$
gives $s^4=1$ for every $s\in S$, hence $P_2(S)=\langle s^4\rangle=1$, while $S'\ne1$;
therefore $S'\not\le P_2(S)$ and $S$ is not powerful by \eqref{eq:powerful}. Since $G$ is a
finite $2$-group whose squares do form a subgroup, the hypotheses of both the first and the
third sentence of 21.137 are met and both are answered in the negative.
\end{proof}

\begin{lemma}\label{lem:vacuous}
In every group $H$ one has $H'\le\langle h^2:h\in H\rangle$.
\end{lemma}

\begin{proof}
Put $K=\langle h^2:h\in H\rangle$. Then $K\trianglelefteq H$, and every element of $H/K$
squares to the identity, so $H/K$ has exponent at most $2$. A group of exponent at most $2$
is abelian, since $xy=(xy)^{-1}=y^{-1}x^{-1}=yx$. Hence $H/K$ is abelian and $H'\le K$.
\end{proof}

\section{Scope}

Only the headline sentence of 21.137 and its third sentence are answered above. The
\emph{second} sentence, ``for $p\ne2$, if the $p$-th powers in a $p$-group of exponent $p^2$
form a subgroup, must that subgroup be abelian?'', is a different question, no object in
this note bears on it, and it is not addressed here. We do not claim that $128$ is the least
order at which the phenomenon occurs, nor that $G$ is the only witness of its order.

We make no priority claim. We have not found Theorem~\ref{thm:main} in print, but our
search of the literature was not exhaustive and $\Syl_2(S_8)$ is an obvious group to test,
so an exercise-level or folklore appearance elsewhere is possible.

\section{Reproduction and verification}

Everything above can be redone from the printed data alone: enter the seven permutations
into any computer algebra system, close under multiplication to obtain $128$ elements, form
the set of squares and test it for closure. Lemma~\ref{lem:structure} and the proof of
Theorem~\ref{thm:main} use nothing beyond \eqref{eq:pres} and Table~\ref{tab:roots} and can
be checked with a pencil.

The accompanying program \texttt{verify.py} re-derives the quantities asserted here in exact
integer and permutation arithmetic, using the Python standard library only and no
group-theory package; \texttt{verify.output.txt} records a run of it. Its input is the
printed data: the seven permutations, the $28$ relators \eqref{eq:pres} and the sixteen
roots of Table~\ref{tab:roots}. From these it verifies the relators, rebuilds $G$ by closure
and confirms $|G|=128$ and $\exp G=8$; checks that the $128$ normal words \eqref{eq:nf} are
distinct and that collection by the printed relators rewrites all $128^2=16384$ products of
normal words correctly, which bounds the presented group by $128$ and so makes
\eqref{eq:pres} a presentation of $G$; tests all $256$ ordered products in $S$ for closure,
obtaining $|S|=16$ and $S=N=G'=\Phi(G)=P_1(G)$; confirms the sixteen square roots of
Table~\ref{tab:roots} and their orders; and computes $\exp S=4$, $P_2(S)=1$ and $|S'|=2$,
whence $S'\not\le P_2(S)$.

\begin{thebibliography}{9}

\bibitem{Kourovka}
E.~I. Khukhro and V.~D. Mazurov (eds.),
\emph{Unsolved Problems in Group Theory. The Kourovka Notebook. No.~21},
Sobolev Institute of Mathematics, Novosibirsk.
Cited from the \LaTeX{} e-print
\href{https://arxiv.org/abs/1401.0300v45}{arXiv:1401.0300v45}; Problem 21.137, proposed by
L.~E. Wilson.

\bibitem{LubotzkyMann}
A.~Lubotzky and A.~Mann,
\emph{Powerful $p$-groups. I. Finite groups},
J. Algebra \textbf{105} (1987), no.~2, 484--505.
\href{https://doi.org/10.1016/0021-8693(87)90211-0}{doi:10.1016/0021-8693(87)90211-0}.
The source of \eqref{eq:powerful}, including the clause that replaces $p$-th powers by
fourth powers when $p=2$.

\bibitem{Wilson2002}
L.~E. Wilson,
\emph{On the power structure of powerful $p$-groups},
J. Group Theory \textbf{5} (2002), no.~2, 129--144.
\href{https://doi.org/10.1515/jgth.5.2.129}{doi:10.1515/jgth.5.2.129}.
Proves, for odd $p$, that in a powerful $p$-group the set of $p^k$-th powers is a subgroup;
21.137 asks for the converse.

\end{thebibliography}

\end{document}
